\documentclass{article}

\textwidth=6.5in
\textheight=8.5in
\voffset=-54pt
\oddsidemargin=0in
\evensidemargin=0in
\setlength{\parskip}{8pt}
\setlength{\parindent}{0pt}

\usepackage{amsmath, amssymb}
\usepackage{ifthen}

\usepackage[inline]{enumitem}
\setlist{topsep=0pt, parsep=8pt}

\usepackage[rgb,table]{xcolor}\convertcolorsUfalse
\usepackage{graphicx}

% change hyperref options
\PassOptionsToPackage{hyphens}{url}
\usepackage[bookmarks,colorlinks,linkcolor=black,citecolor=blue,pdfstartview=FitH,urlcolor=blue]{hyperref}
\hypersetup{pdfpagemode=UseNone}
\usepackage{bookmark}

\usepackage{graphicx}
\usepackage{multicol}
\usepackage{framed}
\usepackage{array}

\usepackage{icomma}

\usepackage{soul}
\usepackage[normalem]{ulem}

\usepackage{pgfplots}
\pgfplotsset{compat=1.18}  
\pgfplotscreateplotcyclelist{mycolor}{black, blue, red, green!70!black, magenta, purple, cyan, orange }  
\pgfplotsset{
  disabledatascaling,
  cycle list=tikzcolor, 
  axis lines=middle,
  every axis plot/.style={no marks, thick},
  every axis/.style={
    enlarge x limits=false,
    enlarge y limits=auto,
    xlabel style={at={(current axis.right of origin)},anchor=west},
    ylabel style={at={(current axis.above origin)},anchor=south},
    % extra x and y ticks for asymptotes
    extra x tick style={grid=major, grid style={dashed,black}},
    extra y tick style={grid=major, grid style={dashed,black}},
    /pgf/number format/1000 sep={},
  },    
  tick label style = {font=\footnotesize, fill=\currentbgcolor, inner sep=0pt},    
  xticklabel shift=1pt,    
  scaled ticks=false,
  scaled y ticks=false,
  unbounded coords=jump,
  samples=101,
  minor tick length=0,
  scatter/classes={
    open={mark=*,draw=black,fill=white, mark size=2, thin},
    closed={mark=*,draw=black,fill=black, mark size=2, thin}
  },
  width=3in,
}
\pgfplotsset{open/.style={only marks, mark=*, fill=white, thin, mark size=2}}
\pgfplotsset{closed/.style={only marks, mark=*, draw=black, fill=black,
    thin, mark size=2}}
\pgfplotsset{labeled/.style={nodes near coords, only marks, mark=*, draw=black, fill=black, thin, mark size=2, point meta=explicit symbolic}}

\pgfplotsset{gridgraph/.style={clip=false, axis line style={shorten
      >=-8pt}, xlabel style={xshift=8pt}, ylabel style={yshift=8pt},
    enlargelimits=false, grid=major, tick style={black}}} 

\usepackage{tikz}
\usetikzlibrary{
  matrix,%
  % enables the snake paths
  decorations.pathmorphing,%
  % enables bigger arrows
  decorations.markings,%
  decorations.pathreplacing,%
  decorations.text,%
  calc,%
  patterns,%
  shapes.geometric,%
  arrows,
  decorations.shapes,
  positioning,
  plotmarks,
  intersections
}
\tikzset{>=latex} % sets default arrow type in tikz
\tikzset{font=\normalsize}
\tikzset{mark size=1.5pt, mark options=thin}
\tikzset{pin distance=4pt,
  every pin edge/.style={<-, thin, shorten <= -2pt}}

\interfootnotelinepenalty=10000
\renewcommand{\thefootnote}{(\arabic{footnote})}

\usepackage{multirow, bigdelim}

% change to \usepackage[sols]{handout} to see solutions
\usepackage[sols]{handout}

\newcommand\iscurrentenumi[1]{%
  \ifthenelse{\equal{#1}{\theenumi}}{}{#1}}
\setenumerate[1]{ref=\#\arabic*}
\setenumerate[2]{ref=\protect\iscurrentenumi{\theenumi}(\alph*)}
\newcommand\enumiiprefix[2]{%
  \ifthenelse{{\equal{#1}{\theenumi}}\AND{\equal{#2}{(\alph{enumii})}}}{}{}%
  \ifthenelse{{\equal{#1}{\theenumi}}\AND{\NOT{\equal{#2}{(\alph{enumii})}}}}{#2}{}%
  \ifthenelse{\equal{#1}{\theenumi}}{}{#1#2}%
}
\setenumerate[3]{ref=\protect\enumiiprefix{\theenumi}{(\alph{enumii})}\roman*}

\newlength{\tipmargin}
\makeatletter

\renewenvironment{framed}{
  \setlength{\tipmargin}{\textwidth-\linewidth}
  \advance\@totalleftmargin-\tipmargin
  \def\FrameCommand{\hspace{\tipmargin}\fbox}%
  \advance\linewidth\tipmargin
  \MakeFramed {\advance\hsize-\width \FrameRestore}%
}
{\endMakeFramed}

\makeatother

\definecolor{purple}{rgb}{0.5,0,1.0}

\usepackage{fancyhdr}
\lhead{\textcolor{white}{This content is protected and may not be shared, uploaded, or distributed.}}
\rhead{}
\rfoot{\vspace*{\fill}\footnotesize\textcopyright{} \emph{Harvard Math Ma}}
\renewcommand{\headrulewidth}{0pt}
\pagestyle{fancy}
\begin{document}

\htitle{The Derivative Function}

\begin{problemonly}
  \begin{framed}

You should be able to:
    \begin{itemize}[topsep=0pt,partopsep=0pt,parsep=8pt,itemsep=0pt]
    \item Use the limit definition to compute the derivative function
      $f'$.

    \item Determine whether a function is \emph{differentiable}, both
      graphically and algebraically, and give examples of functions that
      are differentiable and non-differentiable.

    \item Sketch a rough graph of $f'$ when you have the graph of $f$. 

    \end{itemize}
  \end{framed}
\end{problemonly}

\begin{enumerate}
\item  \emph{Check up!} Suppose $f$ is a function. 

  \begin{enumerate}
  \item \label{derivative-checkup-definition}
    \begin{problem}[\vfill]
      State the definition of $f'(a)$. 
    \end{problem}

    \begin{solution}
      $\displaystyle f'(a) = \lim_{h \to 0} \frac{f(a+h) - f(a)}{h}$
    \end{solution}

  \item
    \begin{problem}[\vfill]
      How do we interpret $f'(a)$ graphically?
    \end{problem}

    \begin{solution}
      $f'(a)$ is the slope of the line tangent to the graph of $y=f(x)$ at $x=a$. 
    \end{solution}

  \item
    \begin{problem}[\vfill]
      If $f(t)$ represents the position at time $t$ of a runner along a
      running path, how could we interpret $f'(a)$? 
    \end{problem}
    
    \begin{solution}
      $f'(a)$ is the runner's instantaneous velocity at time $t=a$. 
    \end{solution}

  \item

    \note{Here, I want students to remember that $\dfrac{f(a + h) -
        f(a)}{h}$ represents an average rate of change / slope of a
      secant line.}

    \begin{problem}[\vfill\newpage]
      Your friend is unsure why your answer to
      \ref{derivative-checkup-definition} represents the derivative.  Give
      your friend a short explanation. 
    \end{problem}

    \begin{solution}
      The expression $\dfrac{f(a+h)-f(a)}{h}$ represents the slope of a secant line between the points $\bigl(a, f(a)\bigr)$
      and $\bigl(a+h, f(a+h)\bigr)$. As we move $x=a+h$ closer to $x=a$, i.e. as $h$ approaches zero, the value of
      this expression approaches the slope
      of the tangent line at $x=a$. We use the word ``limit'' to represent the number that the expression approaches, 
      and we notate it
      as $\displaystyle f'(a) = \lim_{h \to 0} \frac{f(a+h) - f(a)}{h}$.
    \end{solution}
  \end{enumerate}

\item  \label{derivative-function-1}

  Here's the graph of a function $f(t)$; we don't know a formula for $f(t)$, so answer the questions below using just the graph. 
  \begin{center}
    \begin{tikzpicture}
      \begin{axis}[width=2in, xtick={-2, 2}, ytick=\empty]
        \addplot+[domain=-3:3] {x^2 - 4};
      \end{axis}
    \end{tikzpicture}
  \end{center}

  \begin{enumerate}

  \item \label{derivative-function-1-qualitative}

    \note{These questions are encouraging students to start thinking of $f'$ as a function.}

    \begin{enumerate}
    \item
      \begin{problem}[\stretch{0.2}]
        Which is greater, $f'( \frac{1}{2})$ or $f'(1)$?
      \end{problem}

      \begin{solution}
        We can draw tangent lines at $x=\frac{1}{2}$ and $x=1$ and compare their slopes.
        \begin{center}
          \begin{tikzpicture}
            \begin{axis}[width=2in, xtick={-2, 2}, ytick=\empty]
              \addplot+[domain=-3:3] {x^2 - 4};
              \addplot[red, domain=-0.5:1.5] {(x-0.5)-3.75};
              \addplot[blue, domain=0:2] {2*(x-1)-3};
            \end{axis}
          \end{tikzpicture}
        \end{center}

        Both tangent lines have positive slopes, but the slope of the tangent line at $x=1$ is more positive. Thus, 
        $f(1) > f(\frac{1}{2})$. 
      \end{solution}

    \item
      \begin{problem}[\stretch{0.2}]
        Which is greater, $f' (-3)$ or $f'(-1)$?
      \end{problem}

      \begin{solution}
        We can draw tangent lines at $x=-3$ and $x=-1$ and compare their slopes.
        \begin{center}
          \begin{tikzpicture}
            \begin{axis}[width=2in, xtick={-2, 2}, ytick=\empty]
              \addplot+[domain=-3:3] {x^2 - 4};
              \addplot[red, domain=-3:-2] {-6*(x+3)+5};
              \addplot[blue, domain=-2:0] {-2*(x+1)-3};
            \end{axis}
          \end{tikzpicture}
        \end{center}

        Both tangent lines have negative slopes, but the slope of the tangent line at $x=-1$ is less negative. Thus, 
        $f(-1) > f(-3)$. 
      \end{solution}

    \item

      \note{Hopefully students will begin to realize that $f'(t) > 0$ corresponds to $f(t)$ increasing. We will talk about this more next time (and deal with subtleties like: if $f'(t) > 0$, then $f(t)$ is increasing, but if $f(t)$ is increasing and differentiable, we can only say $f'(t) \geq 0$).}

      \begin{problem}[\stretch{0.3}]
        For what values of $t$ will $f'(t)$ be positive?
      \end{problem}

      \begin{solution} 
        Looking at the tangent lines we drew above, we can see that
            $f'(t)>0 $ when $t>0$. 
        We can also see that this coincides with where $f$ is increasing.
      \end{solution}

    \item
      \begin{problem}[\stretch{0.3}]
        For what values of $t$ will $f'(t)$ be negative?
      \end{problem}

      \begin{solution} 
        Similarly, we can see from our sketches that
            $f'(t)<0 $ when $t<0$.
        This coincides with where $f$ is decreasing. 
      \end{solution}

    \item

      \note{For $f'$, I will encourage students to think about how $f'(-a)$ relates to $f'(a)$.}

      \begin{problem}[\stretch{0.5}]
        Is $f$ even, odd, or neither? Is $f'$ even, odd, or neither? Explain.
      \end{problem}

      \begin{solution}
        $f$ is an even function because the graph of $f$ is symmetric over the $y$-axis.
        Since we do not have a graph of $f'$, we should use the algebraic tests for even/odd symmetry. 
        Looking at the graph of $f$, 
        $f'(2)$ is positive and
        $f'(-2)$ is negative. Because of the even symmetry of $f$, it seems reasonable that $f'(-2)=-f'(2)$, and in general,
        $f'(-a)=-f'(a)$. This means that $f'$ is odd.
      \end{solution}

    \end{enumerate}

  \item

    \note{Here, I just want students to sketch an odd increasing function.}

    \begin{problem}[\vfill\newpage]
      Sketch a rough graph of $f'$.
    \end{problem}

    \begin{solution}
      We know the following facts about $f'$.
      \begin{itemize}
        \item $f'$ is positive on $(0, \infty)$.
        \item $f'$ is negative on $(-\infty, 0)$. 
        \item $f'(0)$ is zero.
        \item $f'$ is odd.
      \end{itemize}
      Based on the graph of $f$, we can also deduce that $f'$ will be increasing on $(-\infty, \infty)$. So its graph
      might look something like this.
      \begin{center}
        \begin{tikzpicture}
          \begin{axis}[width=2in, ytick=\empty]
            \addplot[domain=-3:3] {2*x};
          \end{axis}
        \end{tikzpicture}
      \end{center}
    \end{solution}

  \item \label{derivative-function-1-formula}

    \begin{problem}[0.25in]
      One of these is the formula for $f(t)$. Which one?
      \[ t^2 + 4 \qquad t^2 - 4 \qquad (t - 2)^2 \qquad 4t^2 \]
    \end{problem}

    \begin{solution}
      Based on the zeros being $t=\pm 2$, we can deduce that the formula is \fbox{$f(t)=t^2-4$. }
    \end{solution}

  \item
    \begin{problem}[\stretch{1.5}]
      Use the formula you picked in \ref{derivative-function-1-formula} to find a simplified formula for $f'(t)$. Does your result agree with your answers to \ref{derivative-function-1-qualitative}?
    \end{problem}

    \begin{solution} 

      \begin{align*} 
        f'(t) &= \lim_{h \to 0} \frac{ f(t+h) -f(t) }{t+h-t} \\
        &= \lim_{h \to 0} \frac{ \bigl((t+h)^2-4\bigr) - \bigl(t^2-4)\bigr)}{h} \\
        &= \lim_{h \to 0} \frac{ t^2 +2th+h^2 - 4 - t^2 - 4}{h} \\
        &= \lim_{h \to 0} \frac{ 2th +h^2}{h} \\
        &= \lim_{h \to 0} 2t +h \\
        &= \boxed{2t}
      \end{align*} 
      As a way to make sure that this function meets our expectations we note that
      \begin{multicols}{2} \begin{itemize} 
        \item $f'(t) = 2t$ is an increasing function
        \item $f(t)$ is positive when $t>0$ 
        \item $f(t)$ is negative when $t<0$ 
        \item $f(t)$ is zero when $t=0$
        \end{itemize} 
      \end{multicols}

    \end{solution}

  \end{enumerate}

  \probonly{
    \begin{framed}
      \emph{Notation.} The function $f'(t)$ is also written as $\displaystyle \frac{df}{dt}$. If $f(t) = t^2$, as in this problem, we can also write $\displaystyle \frac{d}{dt} t^2$ instead of $f'(t)$. (We read $\dfrac{d}{dt}$ as ``the derivative of''.) 
    \end{framed}
  }

\item  \label{|x-1|-and-derivative}

  \begin{enumerate}
  \item
    \begin{problem}[1.15in]
      Sketch the graph of $|x - 1|$.
    \end{problem}

    \begin{solution}
      This is the graph of $|x|$ shifted right by 1:
      \begin{center}
        \begin{tikzpicture}
          \begin{axis}[xtick={1}, ytick={1}, width=3in, height=2in]
            \addplot+[domain=-3:5] {abs(x-1)};
          \end{axis}
        \end{tikzpicture}
      \end{center}
    \end{solution}

  \item

    \note{ When $x \neq 1$, we can easily see the slope of the tangent line at $x$, so the main question is what happens at $x = 0$. The main intuition we want students to have is that, since $|x - 1|$ isn't locally linear at $x = 1$, it doesn't have a tangent line at $x = 1$. (Students who think of a tangent line as ``a line that just touches the graph'' may be puzzled by this; remind them that, when there is a tangent line, if you zoom in enough on the graph, it should start to look just like the tangent line.)

      It's worthwhile to work through what happens if you try to compute $f'(1)$: that you end up with $\displaystyle \lim_{h \to 0} \frac{|h|}{h}$ and what the graph of $\dfrac{|h|}{h}$ looks like. (Students graphed it in \href{https://people.math.harvard.edu/~jjchen/mathma/docs/workshops/workshop03.html}{Workshop 3}, {{\#1}}.)

      Even after this, many students will be unsure what to draw in the graph of $\frac{d}{dx} | x - 1|$ at $x = 1$, so try to have them reason through it.}

    \begin{problem}[\vfill\newpage]
      Sketch the graph of $\displaystyle \frac{d}{dx} |x - 1|$.   (The
      notation $\displaystyle \frac{d}{dx} |x - 1|$ is shorthand for ``the
      derivative of $|x - 1|$''.)  There's one particularly interesting
      value of $x$ here; what's going on at that value of $x$?
    \end{problem}

    \begin{solution}
      For $x < 1$, the graph has slope $-1$; for $x > 1$, the graph has
      slope 1.  Since $|x - 1|$ is not locally linear at $x = 1$,
      $\frac{d}{dx} |x - 1|$ is undefined there; therefore, this is the
      graph of $\frac{d}{dx} |x - 1|$:
      \begin{center}
        \begin{tikzpicture}
          \begin{axis}[xtick={1}, ytick={-1, 1}, width=3in, height=2in,
            y tick label style={yshift=1pt}]
            \addplot+[domain=-3:1] {-1};
            \addplot+[domain=1:5] {1};
            \addplot+[open] coordinates { (1, -1) (1, 1) };
          \end{axis}
        \end{tikzpicture}
      \end{center}
      We can also look algebraically at $x = 1$. If we write $f(x) =
      |x - 1|$, then
      \begin{align*}
        f'(1) &= \lim_{h \rightarrow 0} \frac{ |1+h-1| - |1-1|}{1+h-1 } \\
        & = \lim_{h \rightarrow 0} \frac{ |h|}{h} 
      \end{align*}
      The graph of $g(h) = \frac{|h|}{h}$ looks like this:
      \begin{center}
      \begin{tikzpicture}
        \begin{axis}[ytick={-1, 1}]
          \addplot+[domain=0:5] [thick] {1};   
          \addplot+[domain=-5:0] [thick] {-1};    
          \addplot[open] []coordinates{(0,-1)};
          \addplot[open] []coordinates{(0,1)};
       	\end{axis}
      \end{tikzpicture}
      \end{center}
      So, we can see that $\displaystyle \lim_{h \rightarrow 0} \frac{
        |h|}{h}$ doesn't exist, so $f'(1)$ doesn't exist. 
    \end{solution}

  \end{enumerate}

\end{enumerate}
\begin{framed}
  \textbf{Definition.} We say a function $f$ is
  \uline{non-differentiable at $c$} if $f'(c)$ does not exist.
  Otherwise, we say that $f$ is \uline{differentiable at $c$}. (For
  both definitions, we only consider $c$ that are in the domain of
  $f$.) 
\end{framed}

\begin{enumerate}[resume]
\item  \label{interpret-derivatives-gromit-copy}

  \note{This is to give students an example of interpreting a derivative that isn't with respect to time.}

  \label{interpret-derivatives-gromit}

  Gromit has been growing a giant squash for Tottington Hall's annual Giant Vegetable Competition. He has carefully tracked his squash's length and weight. Let $w(\ell)$ be the squash's mass in kg when its length is $\ell$ cm.
  \begin{enumerate}
  \item
    \begin{problem}[0.25in]
      What are the units of $w'(\ell)$?  
    \end{problem}

    \begin{solution}
      kg/cm
    \end{solution}

  \item
    \begin{problem}[0.35in]
      Another notation for $w'(\ell)$ is $\dfrac{dw}{d\ell}$. Do you expect $\dfrac{dw}{d\ell}$ to be positive or negative? Why?
    \end{problem}

    \begin{solution}
      As $\ell$ (the squash's length) increases, $w$ (the squash's mass) should increase. That is, $w$ should be an increasing function of $\ell$, so $\dfrac{dw}{d\ell}$ should be positive.
    \end{solution}

  \item

    \note{If we have time, I'll try to emphasize why we have to say ``about 8 kg'', not ``exactly 8 kg''.}

    Suppose $w'(70) = 8$. (This statement can also be written as $\displaystyle \left. \frac{dw}{d\ell} \right|_{\ell=70} = 8$.) Which is the following is the most reasonable conclusion?

    \begin{enumerate}[label=\protect\maybecircle{\Alph*.}]
    \plainitem It takes 70 days for the squash to grow to be 8 kg. 
    \plainitem When the squash is 70 cm long, it weighs 8 kg.
    \plainitem When the squash is 70 cm longer than its current length, it will weigh about 8 kg more than it currently does.
    \circleitem When the squash grows from 70 cm to 71 cm, it will gain about 8 kg in mass.
    \end{enumerate}

  \item

    \note{Make sure students understand the units involved. Students should be able to say something like, ``When the squash is 105 cm long, the instantaneous rate of change of the squash's weight is 10 kg/cm.'' They should also be able to translate this into a more ``everyday'' statement like, ``If the squash grew from 105 cm to 106 cm, we'd expect it to gain about 10 kg.''}

    \begin{problem}[0.35in]
      Interpret the statement $w'(105) = 10$ in words.  
    \end{problem}

    \begin{solution}
      When the squash is 105 cm long, the instantaneous rate of change of the squash's weight is 10 kg/cm. So, for example, if the squash grew from 105 cm to 106 cm, we'd expect it to gain about 10 kg.
    \end{solution}

  \end{enumerate}

  \pspace{\newpage}

\item  \label{sketch-derivative-3}

  \note{Here's a GeoGebra applet: \url{https://www.geogebra.org/m/srjt4zst}}

  Here's the graph of a function $f$.

  \begin{tikzpicture}
    \begin{axis}[gridgraph, ymin=-4, ymax=4, axis equal image,
      xtick={-4, -3, -2, -1, 2, 3, 4}, ytick={-4, -3, ..., 4}]
      \addplot+[domain=-4:{9/11}] {1 / (x - 1) + 3/2}; % 
      \addplot+[domain={6/5}:2] {1 / (x - 1) - 1}; %
      \addplot+[domain=2:4] {1 - (x - 3)^2}; %
      \draw[dashed] (1, -4) -- (1, 4); % 
    \end{axis}
  \end{tikzpicture}

  \begin{enumerate}
  \item

    \note{Since 1 isn't in the domain of $f$, $f$ is non-differentiable at 1. We can understand this graphically (the graph of $f$ doesn't have a meaningful tangent line at $x = 1$) or from the definition (the expression $\dfrac{f(1 + h) - f(1)}{h}$ doesn't make sense because $f(1)$ is undefined).} 

    \begin{problem}[\vfill]
      For what value(s) of $x$ in $(-4, 4)$ is $f$ non-differentiable? 
    \end{problem}

    \begin{solution}
      At $x=1$, because there is vertical asymptote there on the graph of $f$, and at $x=2$, because
      there is a corner there on the graph of $f$.
    \end{solution}

  \item
    \begin{problem}[\vfill]
      For what value(s) of $x$ in $(-4, 4)$ is $f'(x) = 0$?
    \end{problem}

    \begin{solution}
      At $x=3$, because the tangent to the graph of $f$ there is horizontal.
    \end{solution}

  \item
    \begin{problem}[\vfill]
      For what value(s) of $x$ in $(-4, 4)$ is $f'(x) > 0$?
    \end{problem}

    \begin{solution}
      When $2 < x < 3$, which we can also write as $(2, 3)$, because the tangent lines on this interval have positive slopes.
    \end{solution}

  \item
    \begin{problem}[\nstretch{6}]
      Sketch the graph of $f'$ on $(-4, 4)$. 
    \end{problem}

    \begin{solution}
    A good place to start is to use the information we came up in the previous parts:
    \begin{itemize}
      \item $f'(x)$ is undefined when $x=1$ and $x=2$
      \item $f'(x)=0$ when $x=3$
      \item $f'(x)$ is positive when $2 < x < 3$
    \end{itemize}
    We can also deduce, based on the graph of $f$, that $f'$ becomes more and more negative on the interval
    $(-\infty, 1)$ as we go from left to right, so it will be decreasing. On the interval $(1, 2)$, $f'$ is negative
    but increasing (becoming less negative as we go from left to right.) On the interval $(2, 4)$, $f'$ is decreasing and
    changes from positive to negative. This gives us something like the following graph.
    
      \begin{tikzpicture}
        \begin{axis}[gridgraph, ymin=-4, ymax=4, axis equal image, samples=500,
          xtick={-4, -3, -2, -1, 2, 3, 4}, ytick={-4, -3, ..., 4}, restrict y to domain=-4:4]
          \addplot[domain=-4:{9/11}] {1 / (x - 1) + 3/2}; % 
          \addplot[domain={6/5}:2] {1 / (x - 1) - 1}; %
          \addplot[domain=2:4] {1 - (x - 3)^2}; %
          \draw[red, dashed] (1, -4) -- (1, 4); % 
          \addplot[very thick, red, domain=-4:2] {-1 / (x - 1)^2}; % 
          \addplot[very thick, red, domain=2:4] {-2*(x - 3)}; %
          \addplot[red, open] coordinates {(2, -1) (2, 2)};
        \end{axis}
      \end{tikzpicture}      
    \end{solution}
  \end{enumerate}

\item 

  \begin{enumerate}
  \item
    \begin{problem}[\vfill]
      Sketch the graph of $f(x) = (x - 2)^3 + 1$.
    \end{problem}

    \begin{solution}
    
      \begin{tikzpicture}
        \begin{axis}[ytick = \empty]
          \addplot+[domain=-0.1:4.1] {(x-2)^3 + 1}; % 
        \end{axis}
      \end{tikzpicture}      
    \end{solution}

  \item
    \begin{problem}[\vfill]
      Using the graph of $f$, sketch the graph of $f'$. 
    \end{problem}

    \begin{solution}
    
      \begin{tikzpicture}
        \begin{axis}[ytick = \empty]
          \addplot+[domain=-0.1:4.1] {(x-2)^3 + 1}; % 
          \addplot[red, domain=-0.1:4.1] {3*(x-2)^2};
        \end{axis}
      \end{tikzpicture}      
    \end{solution}

  \item
    \begin{problem}[\vfill]
      Sketch the graph of $g(x) = \left| (x - 2)^3 + 1 \right|$.
    \end{problem}

    \begin{solution}
    
      \begin{tikzpicture}
        \begin{axis}[ytick = \empty]
          \addplot+[draw=none, domain=-0.1:4.1] {(x-2)^3 + 1}; % 
          \addplot+[domain=-0.1:4.1] {abs((x-2)^3 + 1)}; % 
        \end{axis}
      \end{tikzpicture}      
    \end{solution}

  \item
    \begin{problem}[\vfill]
      Using the graph of $g$, sketch the graph of $g'$.
    \end{problem}

        \begin{solution}
    
      \begin{tikzpicture}
        \begin{axis}[ytick = \empty]
          \addplot+[draw=none, domain=-0.1:4.1] {(x-2)^3 + 1}; % 
          \addplot+[domain=-0.1:4.1] {abs((x-2)^3 + 1)}; % 
          \addplot[red, domain=-0.1:1] {-3*(x-2)^2};
          \addplot[red, domain=1:4.1] {3*(x-2)^2};
          \addplot[red, open] coordinates {(1, 3) (1, -3)};
        \end{axis}
      \end{tikzpicture}      
    \end{solution}

  \end{enumerate}

\end{enumerate}
\end{document}
